\documentclass[12pt,letterpaper]{article}
\usepackage[margin=0.65in]{geometry}
\usepackage[T1]{fontenc}
\usepackage{lmodern}
\pdfmapfile{+lm.map}
\usepackage{tikz}
\pagestyle{empty}
\setlength{\parindent}{0pt}
\hyphenpenalty=10000\exhyphenpenalty=10000
\begin{document}
\begin{tikzpicture}[remember picture,overlay]\begin{scope}[shift={(current page.north west)},x=1in,y=-1in]
\node[anchor=north west,inner sep=0pt,text width=7.2in,align=left,font=\fontsize{11}{13.75}\selectfont] at (0.65,0.41) {Week 19 / No card inside another / Grades 2--3};
\draw[line width=.35pt] (.65,.70)--(7.85,.70);
\node[anchor=north west,inner sep=0pt,text width=6.7in,align=left,font=\fontsize{9}{11.25}\selectfont] at (0.65,10.43) {Bellingham Math Circle / Week 19 / F19-23-v2};
\node[anchor=north west,inner sep=0pt,text width=0.55in,align=right,font=\fontsize{9}{11.25}\selectfont] at (7.3,10.43) {1};
\node[anchor=north west,inner sep=0pt,text width=7.2in,align=left,font=\fontsize{14}{17.5}\selectfont] at (0.65,0.96) {A card fits inside another when all its pictures are on the other card. The empty card fits inside every card. In a collection, no card may fit inside another. Use a card only once in a collection.};
\node[anchor=north west,inner sep=0pt,text width=7.2in,align=left,font=\fontsize{15}{18.75}\selectfont] at (0.65,2.19) {\textbf{Problem 1:} Find every collection of one or more cards from this deck that follows the rule. How do you know you have them all?};
\draw[line width=.7pt,rounded corners=2pt] (1.83,3.32) rectangle (2.77,4.26);
\draw[line width=1.5pt] (1.9052,3.32) -- (1.9804000000000002,3.32);
\draw[line width=.7pt,rounded corners=2pt] (3.13,3.32) rectangle (4.07,4.26);
\draw[line width=1.5pt] (3.2052,3.32) -- (3.2803999999999998,3.32);
\draw[line width=.95pt] (3.3931999999999998,3.5831999999999997) circle[radius=0.11467999999999999in];
\draw[line width=.95pt,line join=round] (3.8068,3.4513179999999997)--(3.932948,3.680678)--(3.680652,3.680678)--cycle;
\draw[line width=.7pt,rounded corners=2pt] (4.43,3.32) rectangle (5.369999999999999,4.26);
\draw[line width=1.5pt] (4.505199999999999,3.32) -- (4.5804,3.32);
\draw[line width=.95pt] (4.6932,3.5831999999999997) circle[radius=0.11467999999999999in];
\draw[line width=.7pt,rounded corners=2pt] (5.7299999999999995,3.32) rectangle (6.67,4.26);
\draw[line width=1.5pt] (5.805199999999999,3.32) -- (5.8804,3.32);
\draw[line width=.95pt,line join=round] (6.4068,3.4513179999999997)--(6.532947999999999,3.680678)--(6.280652,3.680678)--cycle;
\end{scope}\end{tikzpicture}\mbox{}
\newpage
\begin{tikzpicture}[remember picture,overlay]\begin{scope}[shift={(current page.north west)},x=1in,y=-1in]
\node[anchor=north west,inner sep=0pt,text width=7.2in,align=left,font=\fontsize{11}{13.75}\selectfont] at (0.65,0.41) {Week 19 / No card inside another / Grades 2--3};
\draw[line width=.35pt] (.65,.70)--(7.85,.70);
\node[anchor=north west,inner sep=0pt,text width=6.7in,align=left,font=\fontsize{9}{11.25}\selectfont] at (0.65,10.43) {Bellingham Math Circle / Week 19 / F19-23-v2};
\node[anchor=north west,inner sep=0pt,text width=0.55in,align=right,font=\fontsize{9}{11.25}\selectfont] at (7.3,10.43) {2};
\node[anchor=north west,inner sep=0pt,text width=7.2in,align=left,font=\fontsize{15}{18.75}\selectfont] at (0.65,0.98) {\textbf{Problem 2:} Find the largest collection you can from this deck. Find every collection of that size.};
\draw[line width=.7pt,rounded corners=2pt] (0.79,1.89) rectangle (1.48,2.58);
\draw[line width=1.5pt] (0.8452000000000001,1.89) -- (0.9004000000000001,1.89);
\draw[line width=.95pt] (0.9832000000000001,2.0831999999999997) circle[radius=0.08417999999999999in];
\draw[line width=.95pt] (0.89902,2.30262) rectangle (1.06738,2.47098);
\draw[line width=.7pt,rounded corners=2pt] (1.68,1.89) rectangle (2.37,2.58);
\draw[line width=1.5pt] (1.7351999999999999,1.89) -- (1.7904,1.89);
\draw[line width=.7pt,rounded corners=2pt] (2.57,1.89) rectangle (3.26,2.58);
\draw[line width=1.5pt] (2.6252,1.89) -- (2.6803999999999997,1.89);
\draw[line width=.95pt,line join=round] (3.0667999999999997,1.9863929999999996)--(3.159398,2.154753)--(2.9742019999999996,2.154753)--cycle;
\draw[line width=.7pt,rounded corners=2pt] (3.46,1.89) rectangle (4.15,2.58);
\draw[line width=1.5pt] (3.5152,1.89) -- (3.5704,1.89);
\draw[line width=.95pt] (3.6532,2.0831999999999997) circle[radius=0.08417999999999999in];
\draw[line width=.95pt,line join=round] (3.9568,1.9863929999999996)--(4.049398,2.154753)--(3.8642019999999997,2.154753)--cycle;
\draw[line width=.95pt] (3.56902,2.30262) rectangle (3.73738,2.47098);
\draw[line width=.7pt,rounded corners=2pt] (4.35,1.89) rectangle (5.039999999999999,2.58);
\draw[line width=1.5pt] (4.4052,1.89) -- (4.4604,1.89);
\draw[line width=.95pt] (4.5432,2.0831999999999997) circle[radius=0.08417999999999999in];
\draw[line width=.7pt,rounded corners=2pt] (5.239999999999999,1.89) rectangle (5.93,2.58);
\draw[line width=1.5pt] (5.2951999999999995,1.89) -- (5.3504,1.89);
\draw[line width=.95pt,line join=round] (5.7368,1.9863929999999996)--(5.829397999999999,2.154753)--(5.644202,2.154753)--cycle;
\draw[line width=.95pt] (5.349019999999999,2.30262) rectangle (5.517379999999999,2.47098);
\draw[line width=.7pt,rounded corners=2pt] (6.13,1.89) rectangle (6.82,2.58);
\draw[line width=1.5pt] (6.1852,1.89) -- (6.2404,1.89);
\draw[line width=.95pt] (6.3232,2.0831999999999997) circle[radius=0.08417999999999999in];
\draw[line width=.95pt,line join=round] (6.6268,1.9863929999999996)--(6.719398,2.154753)--(6.5342020000000005,2.154753)--cycle;
\draw[line width=.7pt,rounded corners=2pt] (7.02,1.89) rectangle (7.709999999999999,2.58);
\draw[line width=1.5pt] (7.0752,1.89) -- (7.1304,1.89);
\draw[line width=.95pt] (7.12902,2.30262) rectangle (7.2973799999999995,2.47098);
\draw[black!22,line width=.35pt] (.65,5.05)--(7.85,5.05);
\node[anchor=north west,inner sep=0pt,text width=7.2in,align=left,font=\fontsize{15}{18.75}\selectfont] at (0.65,5.33) {\textbf{Problem 3:} Use the same deck. Keep each starting collection and add as many cards as you can. Does being unable to add a card mean a collection is as large as possible?};
\draw[line width=.7pt,rounded corners=2pt] (0.97,6.4) rectangle (1.75,7.180000000000001);
\draw[line width=1.5pt] (1.0324,6.4) -- (1.0948,6.4);
\draw[line width=.95pt] (1.1884000000000001,6.6184) circle[radius=0.09516in];
\draw[line width=.95pt,line join=round] (1.5316,6.508966)--(1.636276,6.699286)--(1.426924,6.699286)--cycle;
\draw[line width=.95pt] (1.0932400000000002,6.866440000000001) rectangle (1.28356,7.056760000000001);
\draw[line width=.7pt,rounded corners=2pt] (0.97,7.65) rectangle (1.75,8.43);
\draw[line width=1.5pt] (1.0324,7.65) -- (1.0948,7.65);
\draw[line width=.95pt] (1.1884000000000001,7.8684) circle[radius=0.09516in];
\draw[line width=.7pt,rounded corners=2pt] (1.99,7.65) rectangle (2.77,8.43);
\draw[line width=1.5pt] (2.0524,7.65) -- (2.1148,7.65);
\draw[line width=.95pt,line join=round] (2.5516,7.758966)--(2.656276,7.949286)--(2.446924,7.949286)--cycle;
\draw[line width=.95pt] (2.1132400000000002,8.11644) rectangle (2.30356,8.30676);
\draw[line width=.7pt,rounded corners=2pt] (0.97,8.9) rectangle (1.75,9.68);
\draw[line width=1.5pt] (1.0324,8.9) -- (1.0948,8.9);
\draw[line width=.95pt] (1.1884000000000001,9.118400000000001) circle[radius=0.09516in];
\draw[line width=.95pt,line join=round] (1.5316,9.008966000000001)--(1.636276,9.199286)--(1.426924,9.199286)--cycle;
\end{scope}\end{tikzpicture}\mbox{}
\newpage
\begin{tikzpicture}[remember picture,overlay]\begin{scope}[shift={(current page.north west)},x=1in,y=-1in]
\node[anchor=north west,inner sep=0pt,text width=7.2in,align=left,font=\fontsize{11}{13.75}\selectfont] at (0.65,0.41) {Week 19 / No card inside another / Grades 2--3};
\draw[line width=.35pt] (.65,.70)--(7.85,.70);
\node[anchor=north west,inner sep=0pt,text width=6.7in,align=left,font=\fontsize{9}{11.25}\selectfont] at (0.65,10.43) {Bellingham Math Circle / Week 19 / F19-23-v2};
\node[anchor=north west,inner sep=0pt,text width=0.55in,align=right,font=\fontsize{9}{11.25}\selectfont] at (7.3,10.43) {3};
\node[anchor=north west,inner sep=0pt,text width=7.2in,align=left,font=\fontsize{15}{18.75}\selectfont] at (0.65,0.98) {\textbf{Problem 4:} For each deck, arrange every card into exactly one row, with each card fitting inside the next card in its row. A card can be alone in a row. Make as few rows as possible.};
\draw[line width=.7pt,rounded corners=2pt] (1.0,2.06) rectangle (1.62,2.68);
\draw[line width=1.5pt] (1.0496,2.06) -- (1.0992,2.06);
\draw[line width=.95pt] (1.1736,2.2336) circle[radius=0.07564in];
\draw[line width=.7pt,rounded corners=2pt] (2.43,2.06) rectangle (3.0500000000000003,2.68);
\draw[line width=1.5pt] (2.4796,2.06) -- (2.5292000000000003,2.06);
\draw[line width=.95pt] (2.6036,2.2336) circle[radius=0.07564in];
\draw[line width=.95pt,line join=round] (2.8764000000000003,2.146614)--(2.959604,2.297894)--(2.7931960000000005,2.297894)--cycle;
\draw[-stealth,line width=.8pt] (1.74,2.37)--(2.30,2.37);
\node[anchor=north west,inner sep=0pt,text width=4.1in,align=left,font=\fontsize{12}{15.0}\selectfont] at (3.4,2.02) {Example of one nested row: the first card fits inside the next. This is not a whole-deck arrangement.};
\draw[line width=.7pt,rounded corners=2pt] (2.195,3.17) rectangle (3.005,3.98);
\draw[line width=1.5pt] (2.2598,3.17) -- (2.3245999999999998,3.17);
\draw[line width=.7pt,rounded corners=2pt] (3.295,3.17) rectangle (4.105,3.98);
\draw[line width=1.5pt] (3.3598,3.17) -- (3.4246,3.17);
\draw[line width=.95pt] (3.5218,3.3968) circle[radius=0.09882in];
\draw[line width=.95pt,line join=round] (3.8782,3.2831569999999997)--(3.986902,3.480797)--(3.769498,3.480797)--cycle;
\draw[line width=.7pt,rounded corners=2pt] (4.395,3.17) rectangle (5.205,3.98);
\draw[line width=1.5pt] (4.4597999999999995,3.17) -- (4.5245999999999995,3.17);
\draw[line width=.95pt] (4.6217999999999995,3.3968) circle[radius=0.09882in];
\draw[line width=.7pt,rounded corners=2pt] (5.495,3.17) rectangle (6.305,3.98);
\draw[line width=1.5pt] (5.5598,3.17) -- (5.6246,3.17);
\draw[line width=.95pt,line join=round] (6.0782,3.2831569999999997)--(6.186902,3.480797)--(5.969498,3.480797)--cycle;
\draw[black!22,line width=.35pt] (.65,5.14)--(7.85,5.14);
\draw[line width=.7pt,rounded corners=2pt] (0.79,5.48) rectangle (1.48,6.17);
\draw[line width=1.5pt] (0.8452000000000001,5.48) -- (0.9004000000000001,5.48);
\draw[line width=.95pt] (0.9832000000000001,5.6732000000000005) circle[radius=0.08417999999999999in];
\draw[line width=.95pt] (0.89902,5.892620000000001) rectangle (1.06738,6.060980000000001);
\draw[line width=.7pt,rounded corners=2pt] (1.68,5.48) rectangle (2.37,6.17);
\draw[line width=1.5pt] (1.7351999999999999,5.48) -- (1.7904,5.48);
\draw[line width=.7pt,rounded corners=2pt] (2.57,5.48) rectangle (3.26,6.17);
\draw[line width=1.5pt] (2.6252,5.48) -- (2.6803999999999997,5.48);
\draw[line width=.95pt,line join=round] (3.0667999999999997,5.576393)--(3.159398,5.744753)--(2.9742019999999996,5.744753)--cycle;
\draw[line width=.7pt,rounded corners=2pt] (3.46,5.48) rectangle (4.15,6.17);
\draw[line width=1.5pt] (3.5152,5.48) -- (3.5704,5.48);
\draw[line width=.95pt] (3.6532,5.6732000000000005) circle[radius=0.08417999999999999in];
\draw[line width=.95pt,line join=round] (3.9568,5.576393)--(4.049398,5.744753)--(3.8642019999999997,5.744753)--cycle;
\draw[line width=.95pt] (3.56902,5.892620000000001) rectangle (3.73738,6.060980000000001);
\draw[line width=.7pt,rounded corners=2pt] (4.35,5.48) rectangle (5.039999999999999,6.17);
\draw[line width=1.5pt] (4.4052,5.48) -- (4.4604,5.48);
\draw[line width=.95pt] (4.5432,5.6732000000000005) circle[radius=0.08417999999999999in];
\draw[line width=.7pt,rounded corners=2pt] (5.239999999999999,5.48) rectangle (5.93,6.17);
\draw[line width=1.5pt] (5.2951999999999995,5.48) -- (5.3504,5.48);
\draw[line width=.95pt,line join=round] (5.7368,5.576393)--(5.829397999999999,5.744753)--(5.644202,5.744753)--cycle;
\draw[line width=.95pt] (5.349019999999999,5.892620000000001) rectangle (5.517379999999999,6.060980000000001);
\draw[line width=.7pt,rounded corners=2pt] (6.13,5.48) rectangle (6.82,6.17);
\draw[line width=1.5pt] (6.1852,5.48) -- (6.2404,5.48);
\draw[line width=.95pt] (6.3232,5.6732000000000005) circle[radius=0.08417999999999999in];
\draw[line width=.95pt,line join=round] (6.6268,5.576393)--(6.719398,5.744753)--(6.5342020000000005,5.744753)--cycle;
\draw[line width=.7pt,rounded corners=2pt] (7.02,5.48) rectangle (7.709999999999999,6.17);
\draw[line width=1.5pt] (7.0752,5.48) -- (7.1304,5.48);
\draw[line width=.95pt] (7.12902,5.892620000000001) rectangle (7.2973799999999995,6.060980000000001);
\end{scope}\end{tikzpicture}\mbox{}
\newpage
\begin{tikzpicture}[remember picture,overlay]\begin{scope}[shift={(current page.north west)},x=1in,y=-1in]
\node[anchor=north west,inner sep=0pt,text width=7.2in,align=left,font=\fontsize{11}{13.75}\selectfont] at (0.65,0.41) {Week 19 / No card inside another / Grades 2--3};
\draw[line width=.35pt] (.65,.70)--(7.85,.70);
\node[anchor=north west,inner sep=0pt,text width=6.7in,align=left,font=\fontsize{9}{11.25}\selectfont] at (0.65,10.43) {Bellingham Math Circle / Week 19 / F19-23-v2};
\node[anchor=north west,inner sep=0pt,text width=0.55in,align=right,font=\fontsize{9}{11.25}\selectfont] at (7.3,10.43) {4};
\node[anchor=north west,inner sep=0pt,text width=7.2in,align=left,font=\fontsize{15}{18.75}\selectfont] at (0.65,0.98) {\textbf{Problem 5:} Could four cards from this deck follow the collection rule? Explain why your answer covers every choice of four cards.};
\draw[line width=.7pt,rounded corners=2pt] (0.79,2.05) rectangle (1.48,2.7399999999999998);
\draw[line width=1.5pt] (0.8452000000000001,2.05) -- (0.9004000000000001,2.05);
\draw[line width=.95pt] (0.9832000000000001,2.2432) circle[radius=0.08417999999999999in];
\draw[line width=.95pt] (0.89902,2.46262) rectangle (1.06738,2.6309799999999997);
\draw[line width=.7pt,rounded corners=2pt] (1.68,2.05) rectangle (2.37,2.7399999999999998);
\draw[line width=1.5pt] (1.7351999999999999,2.05) -- (1.7904,2.05);
\draw[line width=.7pt,rounded corners=2pt] (2.57,2.05) rectangle (3.26,2.7399999999999998);
\draw[line width=1.5pt] (2.6252,2.05) -- (2.6803999999999997,2.05);
\draw[line width=.95pt,line join=round] (3.0667999999999997,2.1463929999999998)--(3.159398,2.314753)--(2.9742019999999996,2.314753)--cycle;
\draw[line width=.7pt,rounded corners=2pt] (3.46,2.05) rectangle (4.15,2.7399999999999998);
\draw[line width=1.5pt] (3.5152,2.05) -- (3.5704,2.05);
\draw[line width=.95pt] (3.6532,2.2432) circle[radius=0.08417999999999999in];
\draw[line width=.95pt,line join=round] (3.9568,2.1463929999999998)--(4.049398,2.314753)--(3.8642019999999997,2.314753)--cycle;
\draw[line width=.95pt] (3.56902,2.46262) rectangle (3.73738,2.6309799999999997);
\draw[line width=.7pt,rounded corners=2pt] (4.35,2.05) rectangle (5.039999999999999,2.7399999999999998);
\draw[line width=1.5pt] (4.4052,2.05) -- (4.4604,2.05);
\draw[line width=.95pt] (4.5432,2.2432) circle[radius=0.08417999999999999in];
\draw[line width=.7pt,rounded corners=2pt] (5.239999999999999,2.05) rectangle (5.93,2.7399999999999998);
\draw[line width=1.5pt] (5.2951999999999995,2.05) -- (5.3504,2.05);
\draw[line width=.95pt,line join=round] (5.7368,2.1463929999999998)--(5.829397999999999,2.314753)--(5.644202,2.314753)--cycle;
\draw[line width=.95pt] (5.349019999999999,2.46262) rectangle (5.517379999999999,2.6309799999999997);
\draw[line width=.7pt,rounded corners=2pt] (6.13,2.05) rectangle (6.82,2.7399999999999998);
\draw[line width=1.5pt] (6.1852,2.05) -- (6.2404,2.05);
\draw[line width=.95pt] (6.3232,2.2432) circle[radius=0.08417999999999999in];
\draw[line width=.95pt,line join=round] (6.6268,2.1463929999999998)--(6.719398,2.314753)--(6.5342020000000005,2.314753)--cycle;
\draw[line width=.7pt,rounded corners=2pt] (7.02,2.05) rectangle (7.709999999999999,2.7399999999999998);
\draw[line width=1.5pt] (7.0752,2.05) -- (7.1304,2.05);
\draw[line width=.95pt] (7.12902,2.46262) rectangle (7.2973799999999995,2.6309799999999997);
\end{scope}\end{tikzpicture}\mbox{}
\newpage
\begin{tikzpicture}[remember picture,overlay]\begin{scope}[shift={(current page.north west)},x=1in,y=-1in]
\node[anchor=north west,inner sep=0pt,text width=7.2in,align=left,font=\fontsize{11}{13.75}\selectfont] at (0.65,0.41) {Week 19 / No card inside another / Grades 2--3};
\draw[line width=.35pt] (.65,.70)--(7.85,.70);
\node[anchor=north west,inner sep=0pt,text width=6.7in,align=left,font=\fontsize{9}{11.25}\selectfont] at (0.65,10.43) {Bellingham Math Circle / Week 19 / F19-23-v2};
\node[anchor=north west,inner sep=0pt,text width=0.55in,align=right,font=\fontsize{9}{11.25}\selectfont] at (7.3,10.43) {5};
\node[anchor=north west,inner sep=0pt,text width=7.2in,align=left,font=\fontsize{15}{18.75}\selectfont] at (0.65,0.98) {\textbf{Problem 6:} Use this sixteen-card deck. Find the largest collection you can. Explain why you think it is best.};
\draw[line width=.7pt,rounded corners=2pt] (0.79,2.04) rectangle (1.48,2.73);
\draw[line width=1.5pt] (0.8452000000000001,2.04) -- (0.9004000000000001,2.04);
\draw[line width=.95pt] (0.9832000000000001,2.2332) circle[radius=0.08417999999999999in];
\draw[line width=.95pt] (0.89902,2.45262) rectangle (1.06738,2.62098);
\draw[line width=.7pt,rounded corners=2pt] (1.68,2.04) rectangle (2.37,2.73);
\draw[line width=1.5pt] (1.7351999999999999,2.04) -- (1.7904,2.04);
\draw[line width=.95pt,line join=round] (2.1768,2.136393)--(2.2693980000000002,2.3047530000000003)--(2.084202,2.3047530000000003)--cycle;
\draw[line width=.85pt,line join=round] (2.17680,2.42737)--(2.20352,2.50002)--(2.28088,2.50298)--(2.22003,2.55085)--(2.24112,2.62533)--(2.17680,2.58226)--(2.11248,2.62533)--(2.13357,2.55085)--(2.07272,2.50298)--(2.15008,2.50002)--cycle;
\draw[line width=.7pt,rounded corners=2pt] (2.57,2.04) rectangle (3.26,2.73);
\draw[line width=1.5pt] (2.6252,2.04) -- (2.6803999999999997,2.04);
\draw[line width=.7pt,rounded corners=2pt] (3.46,2.04) rectangle (4.15,2.73);
\draw[line width=1.5pt] (3.5152,2.04) -- (3.5704,2.04);
\draw[line width=.95pt] (3.6532,2.2332) circle[radius=0.08417999999999999in];
\draw[line width=.95pt,line join=round] (3.9568,2.136393)--(4.049398,2.3047530000000003)--(3.8642019999999997,2.3047530000000003)--cycle;
\draw[line width=.95pt] (3.56902,2.45262) rectangle (3.73738,2.62098);
\draw[line width=.7pt,rounded corners=2pt] (4.35,2.04) rectangle (5.039999999999999,2.73);
\draw[line width=1.5pt] (4.4052,2.04) -- (4.4604,2.04);
\draw[line width=.95pt] (4.45902,2.45262) rectangle (4.62738,2.62098);
\draw[line width=.85pt,line join=round] (4.84680,2.42737)--(4.87352,2.50002)--(4.95088,2.50298)--(4.89003,2.55085)--(4.91112,2.62533)--(4.84680,2.58226)--(4.78248,2.62533)--(4.80357,2.55085)--(4.74272,2.50298)--(4.82008,2.50002)--cycle;
\draw[line width=.7pt,rounded corners=2pt] (5.239999999999999,2.04) rectangle (5.93,2.73);
\draw[line width=1.5pt] (5.2951999999999995,2.04) -- (5.3504,2.04);
\draw[line width=.95pt] (5.433199999999999,2.2332) circle[radius=0.08417999999999999in];
\draw[line width=.7pt,rounded corners=2pt] (6.13,2.04) rectangle (6.82,2.73);
\draw[line width=1.5pt] (6.1852,2.04) -- (6.2404,2.04);
\draw[line width=.95pt,line join=round] (6.6268,2.136393)--(6.719398,2.3047530000000003)--(6.5342020000000005,2.3047530000000003)--cycle;
\draw[line width=.95pt] (6.23902,2.45262) rectangle (6.40738,2.62098);
\draw[line width=.85pt,line join=round] (6.62680,2.42737)--(6.65352,2.50002)--(6.73088,2.50298)--(6.67003,2.55085)--(6.69112,2.62533)--(6.62680,2.58226)--(6.56248,2.62533)--(6.58357,2.55085)--(6.52272,2.50298)--(6.60008,2.50002)--cycle;
\draw[line width=.7pt,rounded corners=2pt] (7.02,2.04) rectangle (7.709999999999999,2.73);
\draw[line width=1.5pt] (7.0752,2.04) -- (7.1304,2.04);
\draw[line width=.95pt] (7.2132,2.2332) circle[radius=0.08417999999999999in];
\draw[line width=.95pt,line join=round] (7.5168,2.136393)--(7.609398,2.3047530000000003)--(7.424202,2.3047530000000003)--cycle;
\draw[line width=.7pt,rounded corners=2pt] (0.79,2.9299999999999997) rectangle (1.48,3.6199999999999997);
\draw[line width=1.5pt] (0.8452000000000001,2.9299999999999997) -- (0.9004000000000001,2.9299999999999997);
\draw[line width=.85pt,line join=round] (1.28680,3.31737)--(1.31352,3.39002)--(1.39088,3.39298)--(1.33003,3.44085)--(1.35112,3.51533)--(1.28680,3.47226)--(1.22248,3.51533)--(1.24357,3.44085)--(1.18272,3.39298)--(1.26008,3.39002)--cycle;
\draw[line width=.7pt,rounded corners=2pt] (1.68,2.9299999999999997) rectangle (2.37,3.6199999999999997);
\draw[line width=1.5pt] (1.7351999999999999,2.9299999999999997) -- (1.7904,2.9299999999999997);
\draw[line width=.95pt,line join=round] (2.1768,3.0263929999999997)--(2.2693980000000002,3.194753)--(2.084202,3.194753)--cycle;
\draw[line width=.95pt] (1.78902,3.3426199999999997) rectangle (1.95738,3.5109799999999995);
\draw[line width=.7pt,rounded corners=2pt] (2.57,2.9299999999999997) rectangle (3.26,3.6199999999999997);
\draw[line width=1.5pt] (2.6252,2.9299999999999997) -- (2.6803999999999997,2.9299999999999997);
\draw[line width=.95pt] (2.7632,3.1231999999999998) circle[radius=0.08417999999999999in];
\draw[line width=.95pt,line join=round] (3.0667999999999997,3.0263929999999997)--(3.159398,3.194753)--(2.9742019999999996,3.194753)--cycle;
\draw[line width=.95pt] (2.67902,3.3426199999999997) rectangle (2.84738,3.5109799999999995);
\draw[line width=.85pt,line join=round] (3.06680,3.31737)--(3.09352,3.39002)--(3.17088,3.39298)--(3.11003,3.44085)--(3.13112,3.51533)--(3.06680,3.47226)--(3.00248,3.51533)--(3.02357,3.44085)--(2.96272,3.39298)--(3.04008,3.39002)--cycle;
\draw[line width=.7pt,rounded corners=2pt] (3.46,2.9299999999999997) rectangle (4.15,3.6199999999999997);
\draw[line width=1.5pt] (3.5152,2.9299999999999997) -- (3.5704,2.9299999999999997);
\draw[line width=.95pt] (3.6532,3.1231999999999998) circle[radius=0.08417999999999999in];
\draw[line width=.85pt,line join=round] (3.95680,3.31737)--(3.98352,3.39002)--(4.06088,3.39298)--(4.00003,3.44085)--(4.02112,3.51533)--(3.95680,3.47226)--(3.89248,3.51533)--(3.91357,3.44085)--(3.85272,3.39298)--(3.93008,3.39002)--cycle;
\draw[line width=.7pt,rounded corners=2pt] (4.35,2.9299999999999997) rectangle (5.039999999999999,3.6199999999999997);
\draw[line width=1.5pt] (4.4052,2.9299999999999997) -- (4.4604,2.9299999999999997);
\draw[line width=.95pt,line join=round] (4.8468,3.0263929999999997)--(4.939398,3.194753)--(4.754202,3.194753)--cycle;
\draw[line width=.7pt,rounded corners=2pt] (5.239999999999999,2.9299999999999997) rectangle (5.93,3.6199999999999997);
\draw[line width=1.5pt] (5.2951999999999995,2.9299999999999997) -- (5.3504,2.9299999999999997);
\draw[line width=.95pt] (5.433199999999999,3.1231999999999998) circle[radius=0.08417999999999999in];
\draw[line width=.95pt] (5.349019999999999,3.3426199999999997) rectangle (5.517379999999999,3.5109799999999995);
\draw[line width=.85pt,line join=round] (5.73680,3.31737)--(5.76352,3.39002)--(5.84088,3.39298)--(5.78003,3.44085)--(5.80112,3.51533)--(5.73680,3.47226)--(5.67248,3.51533)--(5.69357,3.44085)--(5.63272,3.39298)--(5.71008,3.39002)--cycle;
\draw[line width=.7pt,rounded corners=2pt] (6.13,2.9299999999999997) rectangle (6.82,3.6199999999999997);
\draw[line width=1.5pt] (6.1852,2.9299999999999997) -- (6.2404,2.9299999999999997);
\draw[line width=.95pt] (6.23902,3.3426199999999997) rectangle (6.40738,3.5109799999999995);
\draw[line width=.7pt,rounded corners=2pt] (7.02,2.9299999999999997) rectangle (7.709999999999999,3.6199999999999997);
\draw[line width=1.5pt] (7.0752,2.9299999999999997) -- (7.1304,2.9299999999999997);
\draw[line width=.95pt] (7.2132,3.1231999999999998) circle[radius=0.08417999999999999in];
\draw[line width=.95pt,line join=round] (7.5168,3.0263929999999997)--(7.609398,3.194753)--(7.424202,3.194753)--cycle;
\draw[line width=.85pt,line join=round] (7.51680,3.31737)--(7.54352,3.39002)--(7.62088,3.39298)--(7.56003,3.44085)--(7.58112,3.51533)--(7.51680,3.47226)--(7.45248,3.51533)--(7.47357,3.44085)--(7.41272,3.39298)--(7.49008,3.39002)--cycle;
\end{scope}\end{tikzpicture}\mbox{}
\newpage
\begin{tikzpicture}[remember picture,overlay]\begin{scope}[shift={(current page.north west)},x=1in,y=-1in]
\node[anchor=north west,inner sep=0pt,text width=7.2in,align=left,font=\fontsize{11}{13.75}\selectfont] at (0.65,0.41) {Week 19 / No card inside another / Grades 2--3};
\draw[line width=.35pt] (.65,.70)--(7.85,.70);
\node[anchor=north west,inner sep=0pt,text width=6.7in,align=left,font=\fontsize{9}{11.25}\selectfont] at (0.65,10.43) {Bellingham Math Circle / Week 19 / F19-23-v2};
\node[anchor=north west,inner sep=0pt,text width=0.55in,align=right,font=\fontsize{9}{11.25}\selectfont] at (7.3,10.43) {6};
\node[anchor=north west,inner sep=0pt,text width=7.2in,align=left,font=\fontsize{15}{18.75}\selectfont] at (0.65,0.98) {\textbf{Problem 7:} Arrange all sixteen cards from Problem 6 into as few rows as possible, with each card fitting inside the next card in its row. Every card must appear in exactly one row.};
\draw[black!22,line width=.35pt] (.65,7.05)--(7.85,7.05);
\node[anchor=north west,inner sep=0pt,text width=7.2in,align=left,font=\fontsize{15}{18.75}\selectfont] at (0.65,7.34) {\textbf{Problem 8:} What is the largest possible collection from the sixteen-card deck? Explain why no larger collection can work.};
\end{scope}\end{tikzpicture}\mbox{}
\end{document}
